that is isolated are geometric invariants that are by no means trivial, and will undoubtedly have their part to play in the ``arithmetic'' study of these local rings.

Of course, once one has a reasonable notion of a Chow ring for a regular noetherian scheme admitting an ample invertible Module, one should generalize to this case the relations established or conjectured, respectively, in n\textsuperscript{o} 4 between this ring and the ring @@M0@@, and formulate in terms of this ring the question of the ``Riemann--Roch formula without denominators'' stated in n\textsuperscript{o} 3.

\sgaStableTarget{sga6.e14.bibliography-of-expose-xiv}\subsection*{BIBLIOGRAPHY}
\begin{enumerate}[label={[\arabic*]}]
\item M. F. Atiyah and F. Hirzebruch, The Riemann--Roch theorem for analytic embeddings, Topology, vol. 1, p. 151-166 (1962).
\item A. Borel and J.P. Serre, Le théorème de Riemann--Roch, Bull. Soc. Math. France, vol. 86, p. 97-136 (1958).
\item R. Bott and H. Samelson, Amer. Journ. of Math. vol 80, p. 1004 (1958).
\item C. Chevalley, Les classes d'équivalence rationnelle, I and II, Seminar at the ENS 1958, exposés 2 and 3.
\item C. Chevalley, Classification des groupes de Lie algébriques, Seminar at the ENS 1956/58.
\item A. Dold and D. Puppe, Homologie nicht-additiver Funktoren. Anwendungen, Annales de l'Inst. Fourier t. XI, p. 201-312 (1961).
\item H. Grauert, Ein Theorem der analytischen Garbentheorie, Pub. Math. n\textsuperscript{o} 5, p. 5-64 (1960).
\item A. Grothendieck, Classes de faisceaux et théorème de Riemann--Roch, mimeographed notes prepared by the Institute for Advanced Study, 1957 (work cited as [RRR] in the present exposé, and reproduced in our seminar as an Appendix to Chap. O).
\item A. Grothendieck, La théorie des classes de Chern, Bull. Soc. Math. France vol. 86, p. 137-154 (1958).
